\section{Dissertation Organization}\label{sec:organization}

Chapter 2 establishes the technical and scientific context for the dissertation. Sections 2.1--2.4 review biomimetic and musculoskeletal robots, pneumatic artificial muscles, and models of BPA force--length--pressure behavior. Section 2.5 introduces the human lower-limb biomechanics and OpenSim models used as design benchmarks. Sections 2.6--2.8 then describe spinal CPG architectures, sensory feedback pathways, the Sensory Afferent Database, synthetic nervous systems, and the OpenSim-to-MuJoCo simulation tools that support neuromechanical control. The chapter concludes by explaining how robotic systems can provide repeatable platforms for neuromechanical experiments that would be impractical or unethical to conduct on living subjects.

Chapter 3 describes the experimental, analytical, and computational methods developed to connect actuator behavior to joint performance and neuromechanical control. Sections 3.2 and 3.3 present the BPA force model and the isometric force-characterization experiment. Sections 3.4 and 3.5 describe the pinned-joint and biomimetic four-bar knee test stands and the isometric torque measurements. Sections 3.6--3.8 present the hybrid torque calculation method, the correction terms for constant length offsets, series compliance, and actuator wrapping, and the optimization of actuator routes. Sections 3.9--3.11 document the balance and inverted-pendulum testbeds, the earlier AnimatLab walker, and the OpenSim-to-MuJoCo and SNS-Toolbox workflows used to construct and evaluate the neuromechanical models.

Chapter 4 reports the experimental and computational results. Section 4.1 summarizes the completed foundation studies, including the Sensory Afferent Database and the previously developed modeling tools. Section 4.2 presents the BPA force-characterization data and evaluates the resulting maximum-force and nondimensional force models against the manufacturer's and published models. Section 4.3 reports the isometric knee-torque measurements, identifies the correction terms required to reconcile rigid-body predictions with experimental data, and evaluates their transfer to the biomimetic knee. It also compares the resulting torque envelopes with human benchmarks. Section 4.6 reports the associated follow-up identification and actuator-route redesign. Sections 4.4 and 4.5 document the balance-platform demonstrations and the neuromechanical simulation recordings.

Chapter 5 interprets the results across actuator, transmission, joint, and controller scales. Sections 5.1--5.3 examine the effects of actuator resting length, the nondimensional force formulation, and its relationship to previously published BPA models. Sections 5.5--5.7 evaluate the biomimetic artificial-muscle analogy and the physical sources of error in isometric torque prediction, including series compliance, actuator wrapping, and transmission geometry. Sections 5.4 and 5.8 discuss the implications of the validated models for real-time control and synthesize the force-characterization, knee-torque, sensory-database, and neuromechanical-modeling contributions into a unified actuator-to-controller design methodology.

Chapter 6 presents concrete research opportunities enabled by the models, experiments, and software tools developed in this dissertation. Sections 6.1 and 6.2 explain how researchers can use the OpenSim-to-MuJoCo conversion, SNS-Toolbox coupling, and two-layer CPG implementation on a reproducible musculoskeletal platform. These tools support controlled comparisons of spinal architectures, sensory-feedback pathways, neural lesions, and stimulation protocols. Section 6.3 describes extensions incorporating vestibular, ocular, and cerebellar contributions, while Section 6.4 identifies ways to deploy the framework on embedded hardware for real-time control of physical BPA-driven robots. Section 6.5 extends these tools to gait transitions, reactive balance, and dynamic whole-body maneuvers. These directions are modular applications of the completed work through which future investigators can reuse, compare, and expand the dissertation's validated actuator models, joint models, experimental methods, and neuromechanical-control tools.

Chapter 7 consolidates the dissertation's principal contributions. It summarizes the improved BPA force models and the experimentally corrected prediction of isometric knee torque. It also reviews the validation of the biomimetic knee and optimized actuator routing and the tools developed for neuromechanical simulation and control. The chapter emphasizes the resulting actuator-to-joint design methodology and explains how the completed experimental and computational foundation supports both improved biomimetic robots and controlled investigations of neuromechanical hypotheses.

Appendix A documents the BPA force implementation, muscle-path and moment-arm calculations, compliance model, and wrapping-loss calculation. Appendix B provides the spiking- and non-spiking-neuron equations used by the synthetic nervous system tools and relates those equations to their SNS-Toolbox implementation. Appendix C records the static BPA pathways, joint transformation matrices, bracket reference frames, stiffness and compliance parameters, and adopted correction-term values.
